\section{Related Work}
\label{sec:related-work}

Driven by advances in generative modeling, the landscape of large-scale video models has evolved significantly, particularly in diffusion-based frameworks. Our review focuses on two broad categories: closed-source models and contributions from the open-source community.


\textbf{Closed-source models.}
Closed-source models are mainly those developed by major technology companies, which aims at high-quality, professional video generation due to the extensive resources invested. We organize the notable models released over the past year in chronological order.
In February 2024, OpenAI introduced Sora~\citep{openaisora2024} and marked a significant leap in AI-generated content. In June 2024, Kling~\citep{kuaishou2024kling} by Kuaishou and Luma 1.0~\citep{luma2024dm} by LUMA AI were opened up for public testing, offering powerful video generation capabilities. Meanwhile, Runway unveiled Gen-3~\citep{runway2024gen3}, building upon Gen-2~\citep{runway2024gen2}, to further elevate the standards in video creation. In July 2024, Shengshu AI released Vidu~\citep{bao2024vidu} equipped with self-designed U-ViT~\citep{bao2022all} architecture.
In September 2024, both Kling and Luma undergone updates to version 1.5.
During the same period, MiniMax introduced Hailuo Video~\citep{minimax2024hailuo}, delivering impressive visual results to the public. In October 2024, Pika Labs launched Pika 1.5~\citep{pika2024pika}, enabling users to customize visual and physical attributes in videos. In addition, Meta introduced Movie Gen~\citep{polyak2024moviegencastmedia}, a series of video foundation models that detailed their training processes and applications. In December 2024, Kling advanced to version 1.6 and, concurrently, Google released Veo 2~\citep{veo22025} with an improved understanding of real-world physics and human movement nuances.

These developments underscore the intense global competition in the video generation sector. In this context, our open-source \method has shown competitive or even superior performance over these commercial models across both internal and external benchmarks, leading in multiple attributes.

\textbf{Contributions from open-source community.}
On the other hand, the open-source community has made substantial contributions not only to the holistic video generative models but also to the exploration of essential model components. Diffusion-based video generation models typically build on Stable Diffusion~\citep{rombach2022sd} architecture. It mainly involves three critical modules: an autoencoder to map the original video into a compact latent space, a text encoder to extract text embeddings, and a neural network optimized via the diffusion model to learn the distribution of these video latents.
For network structure, U-Net~\citep{ronneberger2015u}, initially used in image generation~\citep{ho2020denoising,song2020denoising}, is adapted for video generation by incorporating temporal dimensions.
VDM~\citep{ho2022video} extends the 2D U-Net to a 3D version, while another paradigm~\citep{zhou2022magicvideo,wang2023modelscope,guo2023animatediff} introduces 1D temporal attention combined with 2D spatial attention block to reduce computation costs. Notably, Diffusion Transformers (DiT~\citep{dit}), which rely solely on transformer blocks, are superior to U-Net in visual generation tasks~\citep{chen2023pixartalpha}. This structure has also been transferred to video models~\citep{ma2024latte}, giving rise to two common variants: the original DiT~\citep{dit,HaCohen2024LTXVideo} that uses cross-attention for text embeddings and MM-DiT~\citep{genmo2024mochi,kong2024hunyuanvideo}, where text embeddings are concatenated to visual embeddings for full-attention processing.
Regarding autoencoders,
while the early approach~\citep{rombach2022sd} adopts standard VAE~\citep{kingma2013auto}, recent autoencoders like VQ-VAE~\citep{van2017neural} and VQGAN~\citep{esser2020taming} improve model design for better reconstruction and compression. LTX-Video~\citep{HaCohen2024LTXVideo} modifies the VAE decoder to perform the final denoising step and convert latents into pixels, where the missing high-frequency details are generated during decoding. The text encoder is also critical for text-based video generation. Current powerful video generation models primarily utilize the T5 series~\citep{t5} as the main text encoder, often in conjunction with CLIP~\citep{radford2021learning}.
In the case of HunyuanVideo~\citep{kong2024hunyuanvideo}, T5 is replaced with a Multimodal Large Language Model~\citep{llava,li2024mgm} to achieve more robust alignment between text embeddings and visual features.

By integrating these key modules with effective diffusion-based optimization techniques~\citep{ho2020denoising,song2019generative,flow-matching}, multiple promising open-source video generation models have emerged~\citep{genmo2024mochi,kong2024hunyuanvideo,HaCohen2024LTXVideo,opensora,lin2024open,jin2024pyramidal,cogvideox}.
In \method, we have carefully designed or selected each critical component to ensure high-quality video synthesis.
We provide detailed design information and ablation study, which can facilitate the design of future video generation models.


In addition, numerous studies have explored downstream tasks in video generation. These tasks include repainting~\citep{sdinp,propainter}, editing~\citep{sdedit,ip2p,magicbrush,wang2023videocomposer,dreamVideo2}, controllable generation~\citep{controlnet,scedit,wang2024unianimate}, and frame reference generation~\citep{cogvideox,i2vadapter}, often by leveraging adapter-based and ControlNet-like architectures~\citep{controlnet,chu2024diffcomplete} to incorporate user-specified conditions. We also develop various downstream applications based on \method, which have demonstrated excellent performance.



